% GENERATED FILE. Edit canonical lessons, then run npm run generate:books.
% canonical-source-hash: fnv1a64:eea0f7d4266d9e26
% canonical-lessons: MW-C06-hear-answer, MW-W06-uu-matra, MW-W06-ttha, MW-C06-answer, MW-C06-practice

\chapter{Thik Hoon: Answer in Two New Signs}
\label{ch:mw-wellbeing-answer}
\hlchaptermodality{8}

\begin{chapteropening}
\textbf{By the end of this chapter:} \emph{I can answer \mw{हूं} \mw{ठीक} \mw{हूं}, then complete a polite wellbeing exchange in separately scored listening, speaking, reading, and writing tasks.}

A source-attested wellbeing answer whose only new written shapes are isolated before the complete four-skill exchange.
\end{chapteropening}

\section[hū̃ ṭhīk hū̃]{Hear \emph{hū̃ ṭhīk hū̃} — I am fine}

\label{lesson:MW-C06-hear-answer}



\begin{quote}
Write independent \textbf{\mw{ई}}, say that \emph{mhāro} means \textbf{my}, then ask \emph{āp kaiso ho?} once from memory. Keep the page closed.
\end{quote}

\subsection*{What to know first}
\begin{quote}
\emph{hū̃ ṭhīk hū̃} — \textbf{I am fine}
\end{quote}

Hear three groups: \emph{hū̃ | ṭhīk | hū̃}. The first and last groups match. Hold the \emph{ū} and let the ending nasal colour it; do not add a separate final syllable. The middle begins with a curled-back, breathed \emph{ṭh}. Say the line once. Its two unfamiliar written shapes come later.

\subsection*{Your turn}
Hear the question, then answer \emph{hū̃ ṭhīk hū̃}. Switch roles once. Keep this an ear-and-voice exercise.

\begin{tcolorbox}[breakable,colback=teal!4,colframe=teal!35!black,title={Before you move on}]
Say the answer once without looking at any spelling.

Source: \href{https://www.marwaripathshala.com/marwari-lesson-2-english}{Marwari Pathshala, Lesson 2}.
\end{tcolorbox}
\section[ū]{\mw{ू} — hold the vowel below a known sign}

\label{lesson:MW-W06-uu-matra}



\begin{quote}
Write \textbf{\mw{कांई}}, then write \textbf{\mw{म्हारो}} and point to the virāma. Say \emph{hū̃ ṭhīk hū̃}, then write known \textbf{\mw{ह}} once.
\end{quote}

\subsection*{Script — one vowel mark}
\begin{quote}
\textbf{\mw{ू}}
\end{quote}

This is the dependent long-\emph{ū} sign. It hangs below and to the right of a consonant. Add it to known \textbf{\mw{ह}}: \textbf{\mw{ह} + \mw{ू} = \mw{हू}}. Hold the vowel for one beat longer than a short \emph{u}.

\subsection*{Your turn}
Keep \textbf{\mw{ू}} visible. Finger-trace the mark once, copy it once, then add it to \textbf{\mw{ह}} twice. Say \emph{hū} after the hand stops.

\begin{tcolorbox}[breakable,colback=teal!4,colframe=teal!35!black,title={Before you move on}]
Cover the model and write \textbf{\mw{हू}} once.

Source: \href{https://www.unicode.org/charts/PDF/U0900.pdf}{Unicode Devanagari chart}.
\end{tcolorbox}
\section[ṭha]{\mw{ठ} — curl back, then release breath}

\label{lesson:MW-W06-ttha}



\begin{quote}
Write \textbf{\mw{रो}} and point to \textbf{\mw{ो}}. Write \textbf{\mw{हू}}, then say the middle sound group \emph{ṭhīk} without spelling it.
\end{quote}

\subsection*{Script — one consonant}
\begin{quote}
\textbf{\mw{ठ}}
\end{quote}

Read \textbf{\mw{ठ}} as \emph{ṭha}. Curl the tongue tip back, release it, and let a small breath follow. Keep it distinct from dental \textbf{\mw{थ}}, which your hand already knows.

\subsection*{Your turn}
Keep \textbf{\mw{ठ}} visible. Finger-trace once, copy once, then alternate \textbf{\mw{ठ} | \mw{थ}} while saying \emph{ṭha | tha}.

\begin{tcolorbox}[breakable,colback=teal!4,colframe=teal!35!black,title={Before you move on}]
Cover the model and write \textbf{\mw{ठ}} once.

Source: \href{https://www.unicode.org/charts/PDF/U0900.pdf}{Unicode Devanagari chart}.
\end{tcolorbox}
\section[hū̃ ṭhīk hū̃.]{\mw{हूं} \mw{ठीक} \mw{हूं}. — two new signs complete the answer}

\label{lesson:MW-C06-answer}



\begin{quote}
Write \textbf{\mw{म्हारो}} from memory. Say \emph{hū̃ ṭhīk hū̃}, write \textbf{\mw{ू}}, and write \textbf{\mw{ठ}}.
\end{quote}

\subsection*{Script — build three heard groups}
\begin{quote}
\textbf{\mw{हूं} | \mw{ठीक} | \mw{हूं}.}
\end{quote}

Build \textbf{\mw{हूं}} as \textbf{\mw{ह} + \mw{ू} + \mw{ं}}. Build \textbf{\mw{ठीक}} as \textbf{\mw{ठ} + \mw{ी} + \mw{क}}. Every sign except the two from the last lessons was already secure. The repeated first and last group make the line easier to hold in memory.

\subsection*{Your turn}
1. Read \textbf{\mw{हूं} | \mw{ठीक} | \mw{हूं}} once. 2. Cover it for ten seconds. 3. Write all three groups, then say the answer.

\begin{tcolorbox}[breakable,colback=teal!4,colframe=teal!35!black,title={Before you move on}]
Read the answer once from your own handwriting.

Source: \href{https://www.marwaripathshala.com/marwari-lesson-2-english}{Marwari Pathshala, Lesson 2}.
\end{tcolorbox}
\section[Practice]{Practice — ask how someone is and answer}

\label{lesson:MW-C06-practice}



\begin{quote}
Ask and answer one name. Then change the topic: ask how the other person is.
\end{quote}

\subsection*{The exchange}
\begin{quote}
\textbf{\mw{आप} \mw{कैसो} \mw{हो}?}
\end{quote}

>

\begin{quote}
\textbf{\mw{हूं} \mw{ठीक} \mw{हूं}.}
\end{quote}

The first line is polite/formal. The reply reports being fine. Other regions and relationships may choose another pronoun or question form; this chapter assesses the source-attested pair it has taught.

\subsection*{Your turn}
1. \textbf{Listen:} hear one line and identify question or answer. 2. \textbf{Speak:} hear the question and answer without notes. 3. \textbf{Read:} match the printed question to its answer. 4. \textbf{Write:} hear the answer, wait ten seconds, and write it with no model.

Score all four separately. A correct reading does not replace missing writing; a correct written line does not replace the spoken reply.

\begin{tcolorbox}[breakable,colback=teal!4,colframe=teal!35!black,title={Before you move on}]
Stop after one accurate exchange. Speed can wait.

Sources: \href{https://www.marwaripathshala.com/marwari-lesson-2-english}{Marwari Pathshala, Lesson 2} and \href{https://nepal.sil.org/resources/archives/63768}{SIL International, \emph{Marwari–English Dictionary} archive}.
\end{tcolorbox}
